\documentclass[11pt]{article}

\usepackage[margin=1in, top=0.9in, bottom=0.9in]{geometry}
% Recommended, but optional, packages for figures and better typesetting:
\usepackage{microtype}

\usepackage{natbib}
\usepackage[utf8]{inputenc}				% allow utf-8 input
\usepackage[T1]{fontenc}				% use 8-bit T1 fonts
\usepackage{graphicx} % Required for inserting images
\usepackage{hyperref}
\usepackage{amsmath, amssymb, amsthm, mathtools, bm}
\usepackage{thmtools}
\usepackage{thm-restate}
\usepackage[dvipsnames]{xcolor}
\usepackage{caption}
\usepackage{url}
\usepackage{booktabs} % for professional tables
\usepackage{float}
\usepackage[labelformat=simple]{subcaption}
\usepackage{tikz}
\usepackage{nicefrac}					% compact symbols for 1/2, etc.
\usepackage{multirow}

\hypersetup{
	hidelinks,
    colorlinks=true,
    linkcolor=Blue,
    filecolor=Blue,
    citecolor=Blue,
    urlcolor=Blue,
    pdftitle={When LoRA Beats Full Fine-Tuning: A Theoretical Perspective},
}
\newcommand{\ali}[1]{{\color{Green}[{\bfseries A:} #1]}}
\newcommand{\seb}[1]{{\color{Blue}[{\bfseries S:} #1]}}
\newcommand{\rotem}[1]{{\color{red}[{\bfseries R:} #1]}}


%%%%%%%%%%%%%%%%%%%%%%%%%%%%%%%%
% THEOREMS
%%%%%%%%%%%%%%%%%%%%%%%%%%%%%%%%
\theoremstyle{plain}
\newtheorem{theorem}{Theorem}
\newtheorem{proposition}{Proposition}
\newtheorem{lemma}{Lemma}
\newtheorem{corollary}{Corollary}
\theoremstyle{definition}
\newtheorem{definition}{Definition}
\newtheorem{assumption}{Assumption}
\theoremstyle{remark}
\newtheorem{remark}{Remark}


% Attempt to make hyperref and algorithmic work together better:
\newcommand{\theHalgorithm}{\arabic{algorithm}}

\renewcommand\thesubfigure{(\alph{subfigure})}


%mathematical operators
\newcommand{\norm}[1]{\left\lVert#1\right\rVert_2}
\newcommand{\fnorm}[1]{\left\lVert#1\right\rVert_{\mathrm{F}}}
\newcommand{\opnorm}[1]{\left\lVert#1\right\rVert_{\operatorname{op}}}
\newcommand{\inner}[2]{\left\langle #1, #2 \right\rangle}

\DeclareMathOperator*{\argmax}{arg\,max}
\DeclareMathOperator*{\argmin}{arg\,min}


%cursive letters
\newcommand\aaa[0]{\mathcal{A}}
\newcommand\bbb[0]{\mathcal{B}}
\newcommand\ccc[0]{\mathcal{C}}
\newcommand\ddd[0]{\mathcal{D}}
\newcommand\eee[0]{\mathcal{E}}
\newcommand\fff[0]{\mathcal{F}}
\renewcommand\ggg[0]{\mathcal{G}}
\newcommand\hhh[0]{\mathcal{H}}
\newcommand\iii[0]{\mathcal{I}}
\newcommand\jjj[0]{\mathcal{J}}
\newcommand\kkk[0]{\mathcal{K}}
\renewcommand\lll[0]{\mathcal{L}}
\newcommand\mmm[0]{\mathcal{M}}
\newcommand\nnn[0]{\mathcal{N}}
\newcommand\ooo[0]{\mathcal{O}}
\newcommand\ppp[0]{\mathcal{P}}
\newcommand\qqq[0]{\mathcal{Q}}
\newcommand\rrr[0]{\mathcal{R}}
\newcommand\sss[0]{\mathcal{S}}
\newcommand\ttt[0]{\mathcal{T}}
\newcommand\uuu[0]{\mathcal{U}}
\newcommand\vvv[0]{\mathcal{V}}
\newcommand\www[0]{\mathcal{W}}
\newcommand\xxx[0]{\mathcal{X}}
\newcommand\yyy[0]{\mathcal{Y}}
\newcommand\zzz[0]{\mathcal{Z}}


%brackets
\newcommand\rb[1]{\left(#1\right)}
\renewcommand\sb[1]{\left[#1\right]}
\newcommand\cb[1]{\left\{#1\right\}}
\newcommand\supp[1]{supp\left(#1\right)}
% \newcommand\Exp[0]{\mathbb{E}} Please use \E


% Commands by Rotem
\newcommand{\loss}{\mathcal{L}}
\newcommand{\stochasticLoss}{\hat{\mathcal{L}}}
\newcommand{\R}{\mathbb{R}}
\newcommand{\prob}{\mathbb{P}}
\newcommand{\E}{ \mathbb{E} }
\newcommand{\N}{ \mathbb{N} }
\newcommand{\C}{ \mathbb{C} }
\newcommand{\ie}{\emph{i.e.} }
\newcommand{\eg}{\emph{e.g.} }
\newcommand{\bmu}{\boldsymbol{\mu}}
\newcommand{\A}{\boldsymbol{A}}
\newcommand{\CovMat}{\boldsymbol{\Sigma}}
\newcommand{\rank}{\operatorname{rank}}
\newcommand{\trace}{\operatorname{Tr}}
\newcommand{\Var}[0]{\operatorname{\mathbb{V}ar}}
\newcommand{\kernel}{\operatorname{ker}}
\newcommand{\image}{\operatorname{im}}

\newcommand{\bDelta}[0]{\boldsymbol{\Delta}}
\newcommand{\rvx}{{\bf x}}
\newcommand{\rvy}{{\bf y}}
\newcommand{\rvu}{{\bf u}}
\newcommand{\rvv}{{\bf v}}
\newcommand{\rvw}{{\bf w}}

\newcommand{\transpose}{\mathrm{T}}
\newcommand{\zeroVec}{\boldsymbol{0}}
\newcommand{\IdentityMat}{\boldsymbol{I}}


%ALi's commands

%bold vector
\renewcommand\aa[0]{\mathbf{a}}
\newcommand\bb[0]{\mathbf{b}}
\newcommand\cc[0]{\mathbf{c}}
\newcommand\dd[0]{\mathbf{d}}
\newcommand\uu[0]{\mathbf{u}}
\newcommand\pp[0]{\mathbf{p}}
\newcommand\qq[0]{\mathbf{q}}
\newcommand\vv[0]{\mathbf{v}}
\newcommand\xx[0]{\mathbf{x}}
\newcommand\yy[0]{\mathbf{y}}
\newcommand\zz[0]{\mathbf{z}}


%bold matrices
\renewcommand\AA[0]{\mathbf{A}}
\newcommand\BB[0]{\mathbf{B}}
\newcommand\CC[0]{\mathbf{C}}
\newcommand\DD[0]{\mathbf{D}}
\newcommand\XX[0]{\mathbf{X}}
\newcommand\YY[0]{\mathbf{Y}}
\newcommand\II[0]{\mathbf{I}}
\newcommand\UU[0]{\mathbf{U}}
\newcommand\VV[0]{\mathbf{V}}
\newcommand\WW[0]{\mathbf{W}}
\newcommand\ZZ[0]{\mathbf{Z}}
\newcommand\MM[0]{\mathbf{M}}
\newcommand\NN[0]{\mathbf{N}}
\newcommand\QQ[0]{\mathbf{Q}}
\newcommand\PP[0]{\mathbf{P}}
\newcommand\GG[0]{\mathbf{G}}
\newcommand\SSS[0]{\mathbf{S}}
\newcommand\TT[0]{\mathbf{T}}
\newcommand\lam[0]{\mathbf{\Lambda}}
% \newcommand\ee[0]{\mathbb{E}} Please use \E


%paper related new commands
\newcommand{\vecnoise}{\boldsymbol{\varepsilon}} 
\newcommand{\matnoise}{\bm{\xi}} 
\newcommand{\risk}{\mathcal{R}}
\newcommand{\empsol}{\widehat{\AA}}
\newcommand{\fempsol}{\widehat{\BB}}
\newcommand{\lorasol}{\widehat{\AA}_{\mathrm{LoRA}}}
\newcommand{\gdsol}{\widehat{\AA}_{\mathrm{FFT}}}
\newcommand{\initmodel}{\AA_0}
\newcommand{\empcov}{\boldsymbol{\widehat{\Sigma}}}
\newcommand{\trunk}[1]{\operatorname{trunc}_r \left\{ #1 \right\} }
\newcommand{\vectorization}{\operatorname{vec}}
\newcommand{\Cov}{\operatorname{Cov}}
\newcommand\ww[0]{\mathbf{w}}


\newcommand{\circled}[1]{\tikz[baseline=(C.base)]{\node[draw,circle,inner sep=1.5pt](C){#1};}}

% Additional commands for this paper
\newcommand\hh{\mathbf{h}}
\newcommand\ee{\mathbf{e}}
\renewcommand\tt{\mathbf{t}}
\newcommand{\sigmoid}{\sigma}
\newcommand{\spectral}{\rho}
\newcommand{\sel}{\mathcal{S}}

\hypersetup{
    pdftitle={Selectivity Learning Dynamics and Phase Transitions in Gated State Space Models},
}

\title{Selectivity Learning Dynamics and Phase Transitions \\ in Gated State Space Models}
\author{}
\date{}

\begin{document}
\maketitle

\begin{abstract}
State space models with input-dependent gating---\ie selective SSMs---have emerged as powerful sequence models, yet the training dynamics by which selectivity is learned remain poorly understood.
We study this question in the simplest setting that provably requires selectivity: a gated SSM trained on the running parity task.
We prove that linear (non-gated) SSMs cannot compute parity regardless of hidden dimension, while a gated SSM with a single hidden unit suffices.
This establishes that selectivity is both necessary and sufficient.
Experimentally, we observe that the emergence of selectivity during training follows a sharp phase transition: accuracy remains near chance for hundreds of epochs, then jumps to perfect within a narrow window---precisely when the gate learns to discriminate between input values.
\end{abstract}

%═══════════════════════════════════════════════════════════════
\section{Introduction}
\label{sec:intro}
%═══════════════════════════════════════════════════════════════

Sequence modeling is a cornerstone of modern machine learning.
Recurrent neural networks (RNNs) process sequences by maintaining a hidden state that is updated at each timestep, but classical architectures such as the Elman network and LSTM~\citep{hochreiter1997long} have largely been supplanted by Transformers~\citep{vaswani2017attention} due to the difficulty of training through long dependencies.
Recently, structured state space models (SSMs) have emerged as an efficient alternative~\citep{gu2022efficiently}.
In their simplest form, SSMs define a linear recurrence over a hidden state, enabling fast parallel training via convolutions.
However, the key insight of recent work---most notably Mamba~\citep{gu2023mamba}---is that \emph{selectivity} is essential: the state transition must depend on the input so that the model can choose what to remember and what to forget.

This input-dependent gating is not new. The gated recurrent unit (GRU)~\citep{cho2014learning} implements exactly this mechanism: an update gate $\zz_t$, computed as a function of both the current input and the previous state, controls whether the hidden state is updated or preserved.
The GRU can thus be viewed as a discretized SSM with learned selectivity~\citep{gu2023mamba, orvieto2023resurrecting}.

While the importance of selectivity for model \emph{expressiveness} is well appreciated, a fundamental question remains open:
\begin{center}
\emph{How does selectivity emerge during training?}
\end{center}
The existing SSM literature treats selectivity as a static property of the architecture. We study it as a \emph{dynamic} phenomenon: starting from random initialization, the gate must learn to discriminate between inputs. We ask: what does this learning process look like, and what obstacles does gradient-based optimization face?

To isolate the role of selectivity, we study the simplest task that provably requires it: \textbf{running parity}.
Given a binary sequence, the model must output the cumulative parity (XOR) at every position.
This task demands memory---at each step, the model sees a single bit but must know the parity of all preceding bits---and, as we prove, it demands nonlinear gating.

\paragraph{Contributions.}
\begin{itemize}
    \item We prove that no linear SSM, regardless of hidden dimension, can compute running parity (\Cref{thm:linear-impossible}). The proof follows from the observation that linear recurrences produce linear functions of the input, which cannot represent XOR.
    \item We give a constructive proof that a gated SSM with a single hidden unit ($D=1$) suffices (\Cref{thm:gru-sufficient}), establishing that selectivity is both necessary and sufficient.
    \item We formulate a precise experimental framework for studying the training dynamics of selectivity, defining gate selectivity as a measurable order parameter for the phase transition.
\end{itemize}

%═══════════════════════════════════════════════════════════════
\section{Preliminaries}
\label{sec:prelim}
%═══════════════════════════════════════════════════════════════

%───────────────────────────────────────────────────────────────
\subsection{The Running Parity Task}
\label{sec:parity}
%───────────────────────────────────────────────────────────────

\begin{definition}[Running Parity]
\label{def:parity}
Given a binary input sequence $\xx = (x_1, x_2, \ldots, x_T) \in \{0,1\}^T$, the running parity is the sequence $\yy = (y_1, y_2, \ldots, y_T)$ defined by
\begin{equation}
    y_t = \left(\sum_{i=1}^{t} x_i\right) \bmod 2, \quad t = 1, \ldots, T.
\end{equation}
\end{definition}

Equivalently, $y_1 = x_1$ and $y_t = y_{t-1} \oplus x_t$ for $t \geq 2$, where $\oplus$ denotes XOR.
Concretely:
\begin{center}
\begin{tabular}{lcccccc}
\toprule
Position $t$ & 1 & 2 & 3 & 4 & 5 & 6 \\
\midrule
Input $x_t$ & 1 & 0 & 1 & 1 & 0 & 1 \\
Output $y_t$ & 1 & 1 & 0 & 1 & 1 & 0 \\
\bottomrule
\end{tabular}
\end{center}

The critical property of running parity is that it requires \emph{memory}: at time $t$, the model receives only $x_t$ but must know the cumulative count of all past ones. Moreover, the underlying operation is XOR, which is not linearly separable. Both properties make parity a minimal testbed for selective memory.

%───────────────────────────────────────────────────────────────
\subsection{Linear State Space Model}
\label{sec:linear-ssm}
%───────────────────────────────────────────────────────────────

\begin{definition}[Linear SSM]
\label{def:linear-ssm}
A linear state space model with hidden dimension $D$ is defined by the recurrence
\begin{align}
    \hh_t &= \AA\, \hh_{t-1} + \bb\, x_t, \label{eq:linear-state}\\
    \hat{y}_t &= \sigmoid\!\rb{\cc^\transpose \hh_t + d}, \label{eq:linear-output}
\end{align}
where $\AA \in \R^{D \times D}$, $\bb \in \R^D$, $\cc \in \R^D$, $d \in \R$, and $\hh_0 = \zeroVec$.
The sigmoid function $\sigmoid(a) = 1/(1+e^{-a})$ maps the output to a probability.
\end{definition}

The state transition~\eqref{eq:linear-state} is \emph{input-independent} in the sense that $\AA$ and $\bb$ are fixed parameters---the model mixes old memory with new input at the same rate regardless of the input value.
This lack of selectivity has a fundamental consequence, which we formalize in \Cref{sec:theory}.


%───────────────────────────────────────────────────────────────
\subsection{Gated State Space Model}
\label{sec:gated-ssm}
%───────────────────────────────────────────────────────────────

We introduce selectivity through an input-dependent gate, yielding a gated SSM equivalent to the GRU~\citep{cho2014learning}.

\begin{definition}[Gated SSM]
\label{def:gated-ssm}
A gated state space model with hidden dimension $D$ is defined by
\begin{align}
    \zz_t &= \sigmoid\!\rb{\WW_z\, \hh_{t-1} + \uu_z\, x_t + \bb_z}, \label{eq:gate}\\
    \tilde{\hh}_t &= \tanh\!\rb{\WW_h\, \hh_{t-1} + \uu_h\, x_t + \bb_h}, \label{eq:candidate}\\
    \hh_t &= (\mathbf{1} - \zz_t) \odot \hh_{t-1} + \zz_t \odot \tilde{\hh}_t, \label{eq:gated-update}\\
    \hat{y}_t &= \sigmoid\!\rb{\cc^\transpose \hh_t + d}, \label{eq:gated-output}
\end{align}
where $\WW_z, \WW_h \in \R^{D \times D}$ are recurrent weight matrices, $\uu_z, \uu_h \in \R^D$ are input weight vectors, $\bb_z, \bb_h \in \R^D$ are bias vectors, $\cc \in \R^D$ is the output readout vector, $d \in \R$ is the output bias, and $\hh_0 = \zeroVec$.
The operator $\odot$ denotes element-wise multiplication.
\end{definition}

The total number of parameters is $2D^2 + 5D + 1$.

The gate vector $\zz_t \in (0,1)^D$ is the selectivity mechanism:
\begin{itemize}
    \item When $z_{t,i} \approx 0$: the $i$-th dimension of the state is \emph{preserved} ($h_{t,i} \approx h_{t-1,i}$).
    \item When $z_{t,i} \approx 1$: the $i$-th dimension is \emph{overwritten} with the candidate ($h_{t,i} \approx \tilde{h}_{t,i}$).
\end{itemize}
Critically, $\zz_t$ depends on the current input $x_t$ through $\uu_z$.
For parity, the model must learn a gate that opens on 1s and closes on 0s.

\begin{remark}[Connection to Mamba]
The gated update~\eqref{eq:gated-update} can be written as
$\hh_t = \AA_t\, \hh_{t-1} + \BB_t\, \tilde{\hh}_t$
where $\AA_t = \operatorname{diag}(\mathbf{1} - \zz_t)$ and $\BB_t = \operatorname{diag}(\zz_t)$ are input-dependent diagonal matrices. This is precisely the selective state space model formulation of~\citet{gu2023mamba}, where the state transition matrices are functions of the input.
\end{remark}

%═══════════════════════════════════════════════════════════════
\section{Problem Formulation}
\label{sec:problem}
%═══════════════════════════════════════════════════════════════

\paragraph{Training objective.}
We train the gated SSM to minimize the binary cross-entropy loss averaged over all timesteps:
\begin{equation}
    \label{eq:loss}
    \loss = -\frac{1}{T} \sum_{t=1}^{T} \sb{y_t \log \hat{y}_t + (1 - y_t) \log(1 - \hat{y}_t)},
\end{equation}
where $y_t$ is the true running parity and $\hat{y}_t$ is the model prediction.

\paragraph{Dataset.}
For a sequence length $T$, we train on the \emph{complete} dataset of all $2^T$ binary sequences.
There is no train/test split and no sampling noise---the model has access to the full input space.
Any failure to learn is therefore a failure of \emph{optimization}, not of \emph{statistical generalization}.

\paragraph{Optimization.}
We use Adam~\citep{kingma2015adam} with learning rate $\eta = 10^{-3}$, batch size $64$, and \emph{no weight decay}.
Parameters are initialized with Xavier normal initialization for the recurrent matrices $\WW_z, \WW_h$ and normal initialization with standard deviation $0.5$ for the input and output vectors $\uu_z, \uu_h, \cc$.
Biases are initialized to zero.

\paragraph{Gate selectivity.}
To quantify how well the gate discriminates between inputs, we define the \emph{gate selectivity}:
\begin{equation}
    \label{eq:selectivity}
    \sel = \E\sb{\zz_t \mid x_t = 1} - \E\sb{\zz_t \mid x_t = 0},
\end{equation}
where the expectation is over all positions and sequences, averaged across hidden dimensions.
When $\sel \approx 0$, the gate treats all inputs identically (no selectivity).
When $\sel > 0$, the gate opens more for 1s than for 0s, as required for parity.
We show empirically that $\sel$ serves as an \emph{order parameter} for the phase transition: it remains near zero during the accuracy plateau and jumps sharply at the transition.

%═══════════════════════════════════════════════════════════════
\section{Theoretical Results}
\label{sec:theory}
%═══════════════════════════════════════════════════════════════

We establish two results that together show selectivity is both necessary and sufficient for computing parity.

%───────────────────────────────────────────────────────────────
\subsection{Linear SSMs Cannot Compute Parity}
%───────────────────────────────────────────────────────────────

\begin{theorem}
\label{thm:linear-impossible}
For any hidden dimension $D \geq 1$, no linear SSM (\Cref{def:linear-ssm}) can compute the running parity of binary sequences of length $T \geq 2$.
\end{theorem}

\begin{proof}
Unrolling the recurrence~\eqref{eq:linear-state} from $\hh_0 = \zeroVec$ gives
\begin{equation}
    \hh_t = \sum_{k=0}^{t-1} \AA^k \bb\, x_{t-k}.
\end{equation}
Substituting into~\eqref{eq:linear-output}:
\begin{equation}
    \hat{y}_t = \sigmoid\!\rb{\sum_{k=0}^{t-1} w_k\, x_{t-k} + d}, \quad \text{where } w_k = \cc^\transpose \AA^k \bb \in \R.
\end{equation}
Thus $\hat{y}_t$ is the sigmoid of a \emph{linear} function of the inputs $(x_1, \ldots, x_t)$.

At $t = 2$, we require $\hat{y}_2 = \sigmoid(w_0 x_2 + w_1 x_1 + d)$ to agree with $y_2 = x_1 \oplus x_2$ on all four binary inputs.
This requires:
\begin{align}
    \sigmoid(d) < \tfrac{1}{2} \implies d &< 0, \label{eq:xor1}\\
    \sigmoid(w_0 + d) > \tfrac{1}{2} \implies w_0 &> -d, \label{eq:xor2}\\
    \sigmoid(w_1 + d) > \tfrac{1}{2} \implies w_1 &> -d, \label{eq:xor3}\\
    \sigmoid(w_0 + w_1 + d) < \tfrac{1}{2} \implies w_0 + w_1 &< -d. \label{eq:xor4}
\end{align}
From \eqref{eq:xor2} and \eqref{eq:xor3}, $w_0 + w_1 > -2d > -d$, contradicting \eqref{eq:xor4}.
\end{proof}

\begin{remark}
The result holds regardless of the hidden dimension $D$ because increasing $D$ enriches the linear function's coefficients $w_k = \cc^\transpose \AA^k \bb$ but does not introduce nonlinearity. Parity at position $t=2$ is XOR, which is not linearly separable---a fact first observed by~\citet{minsky1969perceptrons}.
\end{remark}


%───────────────────────────────────────────────────────────────
\subsection{A Gated SSM with One Hidden Unit Suffices}
%───────────────────────────────────────────────────────────────

\begin{theorem}
\label{thm:gru-sufficient}
There exists a gated SSM (\Cref{def:gated-ssm}) with hidden dimension $D = 1$ that computes the running parity of binary sequences of arbitrary length.
\end{theorem}

\begin{proof}
We construct explicit parameters. Set $D = 1$ and define:
\begin{equation}
    W_z = 0, \quad U_z = 20, \quad b_z = -10, \quad W_h = -\frac{1}{a}, \quad U_h = 1, \quad b_h = 0,
\end{equation}
where $a = \tanh(1) \approx 0.762$, and set $C = 10$, $d = -10a/2$.
The model operates on two states: $h = 0$ (even parity) and $h = a$ (odd parity).

\emph{Gate behavior.} When $x_t = 0$: $z_t = \sigmoid(-10) \approx 0$, so $h_t \approx h_{t-1}$ (state preserved). When $x_t = 1$: $z_t = \sigmoid(10) \approx 1$, so $h_t \approx \tilde{h}_t$ (state overwritten).

\emph{State flipping.} When the gate opens ($x_t = 1$), the candidate is $\tilde{h}_t = \tanh(-h_{t-1}/a + 1)$:
\begin{itemize}
    \item If $h_{t-1} = 0$ (even): $\tilde{h}_t = \tanh(1) = a$. Transitions to odd.
    \item If $h_{t-1} = a$ (odd): $\tilde{h}_t = \tanh(-1 + 1) = \tanh(0) = 0$. Transitions to even.
\end{itemize}

\emph{Output.} With $C = 10$ and $d = -5a$: $\hat{y}_t = \sigmoid(10 h_t - 5a)$. When $h_t = 0$: $\hat{y}_t = \sigmoid(-5a) \approx 0.02$. When $h_t = a$: $\hat{y}_t = \sigmoid(5a) \approx 0.98$. The predictions are correct for all $t$.

Since the construction uses only the current input $x_t$ and the previous state $h_{t-1}$ (with no dependence on $t$ or $T$), it generalizes to sequences of arbitrary length.
\end{proof}

The contrast between \Cref{thm:linear-impossible} and \Cref{thm:gru-sufficient} is sharp: the \emph{only} difference between the two models is the gate $\zz_t$.
In the linear SSM, the state transition $\AA$ is a fixed matrix; in the gated SSM, the effective transition $\operatorname{diag}(\mathbf{1} - \zz_t)$ is input-dependent.
This single mechanism---selectivity---is what enables parity computation.

\begin{corollary}
\label{cor:selectivity}
Input-dependent gating (selectivity) is both necessary and sufficient for a recurrent model to compute running parity.
\end{corollary}


\paragraph{The representation--optimization gap.}
\Cref{thm:gru-sufficient} shows that $D=1$ is \emph{representationally} sufficient for parity. However, in our experiments, a $D=1$ gated SSM trained with gradient descent never exceeds $87\%$ accuracy: the solution exists but optimization cannot reach it. This gap between what a model can represent and what it can learn via gradient descent---the \emph{representation--optimization gap}---is a central phenomenon in our experimental study.

\vfill

{\small
\bibliographystyle{plainnat}
\bibliography{references}
}

\end{document}
